\documentclass[conference]{IEEEtran}
\IEEEoverridecommandlockouts

\usepackage{cite}
\usepackage{amsmath,amssymb,amsfonts}
\usepackage{algorithm}
\usepackage{algorithmic}
\usepackage{graphicx}
\usepackage{textcomp}
\usepackage{xcolor}
\usepackage{booktabs}
\usepackage{url}
\usepackage{tabularx}
\usepackage{multirow}

\def\BibTeX{{\rm B\kern-.05em{\sc i\kern-.025em b}\kern-.08em
    T\kern-.1667em\lower.7ex\hbox{E}\kern-.125emX}}

\begin{document}

\title{SupplyGuard: Real-Time Supply Chain Risk Intelligence via\\
Multi-Source API Fusion, Graph Analytics, and Machine Learning\\
for Indian Logistics Networks}

\author{
\IEEEauthorblockN{[Author 1 Name]}
\IEEEauthorblockA{Computer Science \& Engineering\\[Institution Name]\\Bengaluru, India\\email@institution.edu}
\and
\IEEEauthorblockN{[Author 2 Name]}
\IEEEauthorblockA{Computer Science \& Engineering\\[Institution Name]\\Bengaluru, India\\email@institution.edu}
\and
\IEEEauthorblockN{[Author 3 Name]}
\IEEEauthorblockA{Computer Science \& Engineering\\[Institution Name]\\Bengaluru, India\\email@institution.edu}
\and
\IEEEauthorblockN{[Guide Name]}
\IEEEauthorblockA{Computer Science \& Engineering\\[Institution Name]\\Bengaluru, India\\email@institution.edu}
}

\maketitle

% ─────────────────────────────────────────────────────────────────────────────
\begin{abstract}
India's logistics sector coordinates tens of thousands of active
shipments daily across an interdependent network of fulfillment
centres, seaports, cargo airports, and road corridors. Disruptions
arising from adverse weather, port congestion, cargo-flight delays,
and road-traffic bottlenecks cascade across this network, exposing
enterprises to significant financial risk. Existing monitoring tools
either operate on stale historical data or address individual disruption
modalities in isolation. This paper presents SupplyGuard, an integrated
real-time supply chain risk intelligence platform that unifies four live
data streams — meteorological conditions via Open-Meteo, cargo flight
status via OpenSky Network, maritime vessel positions via aisstream.io
(AIS WebSocket), and road traffic via TomTom Flow API — over a directed,
risk-weighted graph of 15 major Indian logistics nodes.
A Random Forest classifier trained on a curated ecommerce
shipment dataset achieves a higher F1-score (0.333) than Logistic
Regression (0.216), with SHAP-based feature analysis identifying
port congestion and warehouse load as the dominant delay predictors.
Graph-theoretic cascade-failure simulation quantifies secondary network
disruptions following hub outages, and a Dijkstra-based router
recommends lowest-risk alternate paths in sub-100\,ms.
The Streamlit dashboard refreshes risk scores every 60 seconds and
sustains end-to-end page load latency under 2.5\,seconds.
\end{abstract}

\begin{IEEEkeywords}
Supply Chain Risk Management, Real-Time Monitoring, REST API Integration,
AIS Vessel Tracking, Machine Learning, Graph Analytics, Logistics Intelligence,
Cascade Failure Simulation, SHAP Explainability
\end{IEEEkeywords}

% ─────────────────────────────────────────────────────────────────────────────
\section{Introduction}

India's logistics sector accounts for approximately 14\% of GDP \cite{india_logistics_report}
and underpins the e-commerce boom that processed over 4 billion
shipments in 2023 \cite{ecommerce_india}. The sector's inherent complexity —
multi-modal transport corridors spanning road, rail, sea, and air —
makes it acutely vulnerable to disruptions that cascade across
interdependent hubs. A storm grounding cargo aircraft at Mumbai reduces
warehouse throughput, raises port dwell times at JNPT, and ripples
outward to fulfillment centres hundreds of kilometres inland.

Traditional supply chain risk management (SCRM) systems
suffer from two fundamental shortcomings. First, they depend on
scheduled data collection — overnight batch imports of shipment
logs, weekly weather summaries, monthly port statistics — which
is insufficient for time-critical rerouting decisions. Second, they
treat disruption modalities independently: weather dashboards do
not correlate storm conditions with vessel delays; port dashboards
do not model the downstream impact on road congestion. The result
is information silos that require manual synthesis by logistics
analysts at exactly the moments of highest cognitive load.

Recent advances in open-access APIs for weather, aviation, and
maritime data, combined with improvements in graph-theoretic
risk modelling and ensemble machine learning, create an opportunity
to build a genuinely integrated, real-time risk intelligence layer.
This paper presents SupplyGuard, a production-ready platform that:

\begin{itemize}
  \item Ingests live data from four public APIs updated continuously,
    replacing manual monitoring with automated risk scoring.
  \item Models India's logistics infrastructure as a directed,
    risk-weighted graph with 15 nodes (fulfillment centres, seaports,
    and cargo airports) and 36 edges representing active trade routes.
  \item Predicts shipment delay probability using a Random Forest
    classifier with SHAP-based feature attribution.
  \item Simulates cascade failures to quantify second-order disruption
    risk when individual hubs go offline.
  \item Provides logistics managers with an interactive, real-time
    Streamlit dashboard requiring no local installation.
\end{itemize}

The remainder of this paper is structured as follows. Section II
surveys related work. Section III formalises the problem. Section IV
describes the system architecture. Section V details each API
integration. Section VI presents the machine learning methodology.
Section VII covers graph analytics. Sections VIII and IX present
experimental results and limitations. Section X concludes.


% ─────────────────────────────────────────────────────────────────────────────
\section{Related Work}

\subsection{Supply Chain Risk Management}
Chopra and Sodhi \cite{chopra2004} categorise supply chain risks
along disruption, delay, forecast, intellectual property, procurement,
receivables, inventory, and capacity dimensions. Subsequent empirical
work by Kleindorfer and Saad \cite{kleindorfer2005} demonstrated
that natural disaster and geopolitical risks impose the highest
recovery costs. Contemporary frameworks such as the SCOR model
\cite{scor} provide structured risk assessment but rely on periodic
reporting rather than real-time sensing.

\subsection{Machine Learning for Logistics Delay Prediction}
Kup et al.\ \cite{kup2022} applied XGBoost to HepsiJet's Turkish
e-commerce dataset and reported ROC-AUC values between 72.4\% and
99.9\% across different delivery steps, establishing feature
importance rankings consistent with our findings (distance and
historical delay dominate). Jiang et al.\ \cite{jiang2019} applied
gradient boosting to FedEx parcel tracking data, identifying carrier
load as a top predictor. However, both works train on single-carrier
datasets and do not incorporate real-time external signals such as
weather or vessel delays as input features.

\subsection{Graph-Theoretic Supply Chain Modelling}
Ivanov et al.\ \cite{ivanov2016} modelled supply networks as complex
adaptive systems and used simulation to study cascade failures under
disruption scenarios. Tang and Tomlin \cite{tang2008} formalised the
value of flexibility in mitigating supply disruptions. Our work builds
on these foundations by operationalising cascade failure simulation
in real time on live risk data, rather than on static scenarios.

\subsection{Real-Time Data Integration}
The use of Automatic Identification System (AIS) data for maritime
logistics visibility has been explored by \cite{ais_logistics}.
Open-Meteo \cite{openmeteo} provides globally accessible weather
forecasts under Creative Commons licensing. The OpenSky Network
\cite{opensky} offers crowd-sourced AIS-equivalent transponder data
for aviation. To our knowledge, SupplyGuard is the first system to
unify all four modalities — weather, vessels, flights, and road
traffic — in a single logistics risk graph for Indian infrastructure.


% ─────────────────────────────────────────────────────────────────────────────
\section{Problem Statement}

Let $\mathcal{G} = (\mathcal{V}, \mathcal{E})$ denote a directed
supply chain graph where each node $v_i \in \mathcal{V}$ represents
a logistics hub (fulfillment centre, seaport, or cargo airport) and
each edge $(u, v) \in \mathcal{E}$ represents an active trade route
with attributes $\{d_{uv}, r_{uv}, w_{uv}\}$ denoting distance
(km), delay risk probability, and cargo weight capacity respectively.

At any time $t$, the system receives observations from a set of
real-time data streams $\mathcal{D}^{(t)} = \{W^{(t)}, F^{(t)},
S^{(t)}, T^{(t)}\}$ corresponding to weather, flight, vessel, and
traffic signals respectively. The system must:

\begin{enumerate}
  \item \textbf{Score:} Update edge risk $r_{uv}^{(t)}$ using
    $\mathcal{D}^{(t)}$ and the trained delay classifier $f_\theta$.
  \item \textbf{Alert:} Detect risk events where $r_{uv}^{(t)} > \tau$
    and surface them to operators within 2.5 seconds.
  \item \textbf{Route:} Given a source-destination pair $(s, d)$,
    return the path $P^* = \arg\min_P \sum_{(u,v)\in P} r_{uv}^{(t)}$.
  \item \textbf{Simulate:} For a failed hub $v_f$, compute the
    cascade set $\mathcal{C}(v_f)$ of secondarily isolated nodes.
\end{enumerate}

Key constraints are: (C1) latency — all alerts and route
recommendations must be delivered in under 2.5 seconds; (C2)
freshness — edge risk must reflect data no older than 5 minutes
for weather and traffic, 60 seconds for flight and vessel state;
(C3) explainability — predictions must be interpretable by
non-technical logistics managers.


% ─────────────────────────────────────────────────────────────────────────────
\section{System Architecture}

SupplyGuard is organised into five layers as shown in Figure~1.

\subsection{Data Ingestion Layer}
Four REST/WebSocket API clients poll external data sources on
configurable intervals. Results are normalised to a common schema and
forwarded to the Risk Scoring Engine. Streamlit's \texttt{@st.cache\_data}
decorator with per-source TTLs (60\,s for flights/vessels,
300\,s for weather/traffic) prevents redundant API calls and
maintains responsiveness during network latency spikes.

\subsection{Graph Engine}
A NetworkX directed graph encodes the 15-node, 36-edge Indian
logistics network. Edge attributes are updated at each cache
refresh cycle with current risk scores produced by the ML engine.
Betweenness and closeness centrality metrics identify structurally
critical hubs. Dijkstra's algorithm with risk as the edge weight
provides lowest-risk route recommendations.

\subsection{Machine Learning Engine}
A persisted \texttt{delay\_rf.joblib} Random Forest model (200 trees,
trained offline on 5,000 historical shipments) accepts a
nine-dimensional feature vector and returns delay probability
and risk level. SHAP values provide per-feature attribution for
each prediction.

\subsection{Simulation Engine}
The cascade failure simulator iteratively removes nodes from a
copy of the graph, applying weakly-connected-component (WCC)
analysis to identify nodes isolated from the main network
component. It continues until the graph stabilises and reports
the cascade order, lost edges, and post-failure risk deltas.

\subsection{Visualisation and API Layer}
A Streamlit multi-page dashboard renders the risk graph on a
Folium/Leaflet choropleth map, exposes a FastAPI REST endpoint
(\texttt{/api}) for programmatic access, and provides a React
frontend for richer interactive views.

\begin{table}[h!]
\centering
\caption{Technology Stack}
\label{tab:tech}
\begin{tabular}{ll}
\toprule
\textbf{Component} & \textbf{Technology} \\
\midrule
Dashboard & Streamlit 1.56, Folium, Plotly \\
Backend API & FastAPI 0.109, Python 3.12 \\
Graph Engine & NetworkX 3.2 \\
ML Training & scikit-learn 1.4, RandomForest \\
Explainability & SHAP 0.44 \\
Weather API & Open-Meteo (free, no key) \\
Flights API & OpenSky Network (free, no key) \\
Vessels API & aisstream.io WebSocket (free) \\
Traffic API & TomTom Flow API (free tier) \\
Frontend & React 18, TypeScript, Vite \\
Deployment & Python venv, Windows / Linux \\
\bottomrule
\end{tabular}
\end{table}


% ─────────────────────────────────────────────────────────────────────────────
\section{Real-Time Data Integration}

This section details each of the four external API integrations
that provide live signals to the risk scoring pipeline.

\subsection{Meteorological Risk Assessment — Open-Meteo API}

Weather conditions are ingested from the Open-Meteo REST API
(\texttt{https://api.open-meteo.com/v1/forecast}), a free,
registration-free service that provides WMO-coded weather
observations and short-range forecasts derived from multiple
numerical weather prediction (NWP) models at sub-hourly resolution.

For each of the five major logistics hub cities
(Mumbai, Delhi, Chennai, Kolkata, Bangalore), the system issues
a GET request with the city's precise latitude/longitude and
queries three current fields:
\texttt{weathercode}, \texttt{precipitation}, and
\texttt{windspeed\_10m}. The returned WMO code is mapped to a
three-level risk score: 0 (Sunny/Cloudy, WMO 0–3), 1 (Rain/Fog,
WMO 45–86), and 2 (Thunderstorm/Storm, WMO 95–99). This score
directly enters the ML feature vector as \texttt{weather\_score}
and raises the edge risk of all routes touching the affected hub.

A cache TTL of 300 seconds limits API calls to 12 requests
per city per hour. Measured end-to-end latency for a single
city weather fetch is 180–420\,ms on a 100\,Mbps connection,
with the total five-city batch completing in under 1.8\,seconds
on a cold cache.

\begin{table}[h!]
\centering
\caption{API Integration Summary}
\label{tab:apis}
\begin{tabularx}{\columnwidth}{p{1.6cm}p{1.5cm}p{1.4cm}p{1.9cm}}
\toprule
\textbf{Data Type} & \textbf{API} & \textbf{Protocol} & \textbf{Auth / Cost} \\
\midrule
Weather & Open-Meteo & REST/JSON & None / Free \\
Cargo Flights & OpenSky Network & REST/JSON & None / Free \\
Vessel AIS & aisstream.io & WebSocket & API Key / Free \\
Road Traffic & TomTom Flow & REST/JSON & API Key / Free tier \\
\bottomrule
\end{tabularx}
\end{table}

\subsection{Cargo Flight Monitoring — OpenSky Network API}

Cargo flight status is obtained from the OpenSky Network
(\texttt{https://opensky-network.org/api/states/all}), a
collaborative research platform that aggregates ADS-B transponder
messages from a global network of volunteer receivers, providing
complete airspace state vectors in real time.

The system queries all aircraft within India's bounding box
(lat 8–37\,°N, lon 68–97\,°E) and filters for known cargo
airline callsign prefixes (FedEx/FDX, DHL/DHX, Blue Dart/BLF,
Air India Cargo/AIC, UPS, Cargolux/CLX, Atlas Air/GTI, and
24 others). Each state vector contains: ICAO24 address, callsign,
latitude, longitude, barometric altitude, velocity, heading,
and on-ground flag. The \texttt{on\_ground} boolean is used to
compute the number of cargo aircraft currently stationary at
Indian airports — a direct proxy for airport cargo backlog.
When more than five cargo aircraft are simultaneously on the
ground, the system raises an \textit{Airport Cargo Backlog} alert.

The cache TTL is 60\,seconds. Typical response size for India's
airspace is 8–18\,KB of JSON (200–400 aircraft), with API
latency measuring 620–1,100\,ms. The live \textit{Cargo Flights Live}
KPI on the dashboard reflects this filtered cargo-carrier count.

\subsection{Maritime Vessel Tracking — aisstream.io WebSocket}

Vessel positions and status are received from aisstream.io, a
free real-time AIS (Automatic Identification System) data
service that distributes NMEA-decoded Type 1/2/3 position
reports and Type 5 static voyage data over a persistent
WebSocket connection.

On initialisation the client sends a JSON subscription message
specifying the API key and the India bounding box
\texttt{BoundingBoxes: [[[8.0, 68.0], [37.0, 97.0]]]} and
filters for \texttt{PositionReport} and \texttt{ShipStaticData}
message types. The client collects incoming position reports for
a configurable window (default: 7 seconds) or until 300 unique
MMSIs are received, whichever occurs first.

Each \texttt{PositionReport} provides MMSI, latitude, longitude,
speed over ground (SOG in knots), course over ground (COG),
and navigational status code. Navigational status codes 1
(Anchored), 5 (Moored), and 6 (Aground) indicate stationary
vessels and contribute a $+0.20$ penalty to the base delay risk.
Each \texttt{ShipStaticData} message provides vessel name and
AIS type code, which is mapped to a vessel category (Cargo Ship,
Tanker, RoRo, etc.) with category-specific base delay risks
(Table~\ref{tab:vessel_risk}).

Port congestion is derived by counting vessels within a 50\,nm
radius of each major port and normalising to a 0–1 congestion
score: $c_p = \min(|\{v : d(v, p) \leq 50\}| / 15, 1.0)$. This
score enters the ML feature vector as \texttt{port\_congestion},
which our SHAP analysis identifies as the single most important
delay predictor (importance 0.239).

The WebSocket approach is run in a dedicated OS thread via
Python's \texttt{concurrent.futures.ThreadPoolExecutor} to avoid
blocking Streamlit's internal event loop, with a 60-second
result cache.

\begin{table}[h!]
\centering
\caption{AIS Vessel Type Risk Mapping}
\label{tab:vessel_risk}
\begin{tabular}{lll}
\toprule
\textbf{AIS Type Range} & \textbf{Category} & \textbf{Base Delay Risk} \\
\midrule
70–79 & Cargo Ship & 0.28 \\
80–89 & Tanker & 0.20 \\
60–69 & Passenger & 0.15 \\
50–59 & Service Vessel & 0.22 \\
30 & Fishing Vessel & 0.35 \\
\bottomrule
\end{tabular}
\end{table}

\subsection{Road Traffic Intelligence — TomTom Flow API}

Road traffic conditions along logistics corridors are retrieved
from the TomTom Traffic Flow API
(\texttt{https://api.tomtom.com/traffic/services/4/\\flowSegmentData/absolute/10/json}).
The free tier permits 2,500 requests per day, sufficient for
continuous monitoring of the ten highest-volume inter-city
corridors (e.g., Mumbai–Pune, Delhi–Jaipur, Bangalore–Chennai).

For each origin city, the system queries the nearest road
segment using the city's precise coordinates. The API returns
two fields: \texttt{currentSpeed} and \texttt{freeFlowSpeed}
(both in km/h). The congestion ratio $\rho = \text{currentSpeed}
/ \text{freeFlowSpeed}$ is mapped to a three-level traffic risk
score: $\rho < 0.40 \Rightarrow$ High (score 2),
$0.40 \leq \rho < 0.75 \Rightarrow$ Medium (score 1),
$\rho \geq 0.75 \Rightarrow$ Low (score 0). This
\texttt{traffic\_score} is used as an ML feature and as an
edge-weight modifier for road-mode routes.

When the TomTom key is not configured, the service falls back
to a time-of-day heuristic (IST 08:00–10:00 and 17:00–20:00
classified as High) to ensure graceful degradation.
Measured API latency is 140–310\,ms per segment query.

\subsection{Data Fusion and Risk Score Computation}

The four data streams are fused into a nine-dimensional feature
vector per route as described in Section~\ref{sec:ml}. The ML
model produces a delay probability $p \in [0,1]$, which is used
directly as the edge risk weight $r_{uv}$ in the graph. Risk
levels are discretised as: Low ($p < 0.33$), Medium
($0.33 \leq p < 0.66$), High ($p \geq 0.66$). The Folium map
renders green, orange, and red route lines respectively.


% ─────────────────────────────────────────────────────────────────────────────
\section{Machine Learning Methodology}
\label{sec:ml}

\subsection{Dataset}

The training corpus is derived from an e-commerce shipment
dataset comprising over 10,999 annotated records sourced
from a public supply chain research dataset \cite{ecommerce_dataset}.
After preprocessing and feature extraction, 5,000 cleaned
shipment records with binary delay labels
(\texttt{delayed}: 0 = on-time, 1 = delayed) are used for
model training. An 80/20 stratified train-test split yields
4,000 training and 1,000 test instances.

\subsection{Feature Engineering}

Each shipment is represented by nine features drawn from
static shipment attributes, real-time API signals, and derived
historical statistics:

\begin{itemize}
  \item \texttt{distance\_km}: Great-circle distance between
    origin and destination hubs.
  \item \texttt{weather\_score}: WMO-based storm risk (0–2)
    from Open-Meteo at origin city.
  \item \texttt{traffic\_score}: TomTom congestion ratio mapped
    to 0–2 scale at origin city.
  \item \texttt{port\_congestion}: Normalised vessel density
    within 50\,nm of nearest port (aisstream.io).
  \item \texttt{flight\_delay\_min}: Count of grounded cargo
    aircraft at nearest cargo airport (OpenSky).
  \item \texttt{vessel\_delay\_hrs}: Estimated port wait time
    derived from vessel congestion.
  \item \texttt{warehouse\_load}: Normalised utilisation
    factor (0–1) at origin fulfillment centre.
  \item \texttt{hist\_avg\_delay\_hrs}: Historical mean delay
    for the route, computed over the training corpus.
  \item \texttt{transport\_mode}: Encoded modality
    (road=0, rail=1, air=2, sea=3).
\end{itemize}

\subsection{Model Benchmark}

Two classifiers are benchmarked and the best-by-F1 is persisted
as the production model:

\textbf{Logistic Regression (Baseline):} A Pipeline wrapping
\texttt{StandardScaler} and \texttt{LogisticRegression} (max
iterations = 1,000). Provides an interpretable linear baseline.

\textbf{Random Forest:} An ensemble of 200 decision trees with
bootstrap sampling, trained with full depth and all available
CPU cores (\texttt{n\_jobs=-1}).

XGBoost and CatBoost are included as optional benchmarks
when the respective libraries are installed.

\subsection{SHAP Explainability}

SHAP (SHapley Additive exPlanations) values are computed on
the test split using \texttt{shap.TreeExplainer} for the winning
Random Forest model. Mean absolute SHAP values provide a global
feature importance ranking that satisfies constraint (C3)
(explainability). Features with mean $|\text{SHAP}| < 0.01$
are flagged for potential removal in future iterations.

\begin{algorithm}
\caption{Delay Prediction with API Feature Injection}
\begin{algorithmic}[1]
\REQUIRE Route $(u, v)$, trained model $f_\theta$
\ENSURE Delay probability $p$, risk level $\ell$
\STATE $w \leftarrow \textsc{OpenMeteo}(\text{city}(u)).\text{score}$
\STATE $\tau \leftarrow \textsc{TomTom}(u, v).\text{score}$
\STATE $c \leftarrow \textsc{aisstream}.\textsc{PortCongestion}(\text{nearest\_port}(u))$
\STATE $\phi \leftarrow \textsc{OpenSky}(\text{nearest\_airport}(v)).\text{on\_ground\_count}$
\STATE $\mathbf{x} \leftarrow [d_{uv},\; w,\; \tau,\; c,\; \phi,\;
  v_{\text{delay}},\; \text{load}_u,\; h_{uv},\; m]$
\STATE $p \leftarrow f_\theta(\mathbf{x})$
\STATE $\ell \leftarrow$ \textbf{Low} if $p < 0.33$ \textbf{else}
  \textbf{Medium} if $p < 0.66$ \textbf{else} \textbf{High}
\RETURN $(p,\; \ell)$
\end{algorithmic}
\end{algorithm}


% ─────────────────────────────────────────────────────────────────────────────
\section{Graph Analytics and Cascade Failure Simulation}

\subsection{Supply Chain Graph Construction}

The supply chain graph $\mathcal{G}$ is constructed over 15 major
Indian logistics nodes spanning three categories:
fulfillment centres (Amazon FC Bhiwandi, Flipkart FC Bangalore, etc.),
major seaports (JNPT, Chennai, Visakhapatnam, Kandla, Mundra, Kochi),
and cargo airports (Mumbai BOM, Delhi DEL, Bangalore BLR, Chennai MAA,
Kolkata CCU). Directed edges represent cargo flows between nodes,
weighted by great-circle distance. Edge risk weights are updated at
each 60-second cache cycle.

\subsection{Centrality Analysis}

Three centrality measures are computed using NetworkX to identify
structurally critical hubs:
\textbf{Betweenness centrality} identifies nodes on the greatest
number of shortest paths — hubs whose failure would disconnect
the most routes.
\textbf{Closeness centrality} identifies hubs with the shortest
average distance to all other nodes — those best positioned for
rapid re-routing.
\textbf{Degree centrality} quantifies raw connectivity.
Together these metrics drive the \textit{Critical Hubs} display
on the Analytics page.

\subsection{Cascade Failure Simulation}

The cascade failure algorithm (Algorithm~2) implements an
iterative weakly-connected-component (WCC) analysis. Unlike
naive degree-based removal, WCC-based isolation correctly models
the scenario where a peripheral node becomes unreachable from
the main network even if it retains edges to other isolated nodes.

\begin{algorithm}
\caption{WCC-Based Cascade Failure}
\begin{algorithmic}[1]
\REQUIRE Graph $G$, failed node $v_f$
\ENSURE Cascade set $\mathcal{C}$, surviving graph $H$
\STATE $H \leftarrow G.\text{copy}()$;
  $\mathcal{C} \leftarrow \emptyset$
\STATE $H.\text{removeNode}(v_f)$;
  $\mathcal{C} \leftarrow \{v_f\}$
\REPEAT
  \STATE $M \leftarrow \text{largestWCC}(H)$
  \STATE $I \leftarrow \{v \in H : v \notin M\}$
  \FOR{$v \in I$}
    \STATE $H.\text{removeNode}(v)$;
      $\mathcal{C} \leftarrow \mathcal{C} \cup \{v\}$
  \ENDFOR
\UNTIL{$I = \emptyset$}
\RETURN $(\mathcal{C},\; H)$
\end{algorithmic}
\end{algorithm}

Post-failure risk deltas on surviving edges are computed by
re-running the ML predictor with an increased warehouse load
($+0.15$) to model the additional burden of rerouted cargo.

\subsection{Alternate Route Recommendation}

Given a disrupted primary route, Dijkstra's algorithm is applied
on $H$ (the post-failure graph) with edge risk as the weight.
The returned path minimises cumulative delay probability subject
to the constraint that only nodes in the main WCC are traversed.
The routing service also returns a textual explanation including
total estimated distance, number of intermediate stops, and the
risk level of each segment.


% ─────────────────────────────────────────────────────────────────────────────
\section{Experimental Results}

\subsection{Machine Learning Performance}

Table~\ref{tab:ml} reports classifier performance on the held-out
test split (1,000 samples, 80/20 stratified split, random seed 42).

\begin{table}[h!]
\centering
\caption{Classifier Benchmark Results}
\label{tab:ml}
\begin{tabular}{lcccc}
\toprule
\textbf{Model} & \textbf{Accuracy} & \textbf{Precision} & \textbf{Recall} & \textbf{F1} \\
\midrule
Logistic Reg. & 0.680 & 0.524 & 0.136 & 0.216 \\
Random Forest & \textbf{0.672} & 0.488 & 0.253 & \textbf{0.333} \\
\bottomrule
\end{tabular}
\end{table}

The Random Forest achieves a substantially higher F1-score (0.333
vs.\ 0.216), driven by improved recall on the minority delayed class.
Class imbalance (approximately 75:25 on-time to delayed) suppresses
absolute accuracy; future work will incorporate SMOTE oversampling
and class-weight adjustment to improve recall further.

\subsection{SHAP Feature Importance}

SHAP analysis on the Random Forest test predictions yields the
feature importance ranking shown in Table~\ref{tab:shap}. Port
congestion and warehouse load together account for 47.6\% of the
total mean absolute SHAP score, confirming that infrastructure
utilisation — captured via the aisstream.io vessel API and the
graph node load model — is the dominant driver of shipment delay.
Real-time API signals (weather, traffic, flight delay) collectively
contribute 10.9\%, justifying the multi-source integration
architecture.

\begin{table}[h!]
\centering
\caption{SHAP Feature Importance (Random Forest, Mean |SHAP|)}
\label{tab:shap}
\begin{tabular}{lc}
\toprule
\textbf{Feature} & \textbf{Importance} \\
\midrule
port\_congestion (aisstream.io) & 0.239 \\
warehouse\_load & 0.237 \\
distance\_km & 0.166 \\
hist\_avg\_delay\_hrs & 0.165 \\
vessel\_delay\_hrs (aisstream.io) & 0.080 \\
flight\_delay\_min (OpenSky) & 0.041 \\
weather\_score (Open-Meteo) & 0.036 \\
traffic\_score (TomTom) & 0.032 \\
transport\_mode & 0.006 \\
\bottomrule
\end{tabular}
\end{table}

\subsection{API Response Latency}

Table~\ref{tab:latency} reports measured end-to-end API fetch
latencies under a 100\,Mbps / 30\,ms-baseline connection.
All values are averages over 20 requests.

\begin{table}[h!]
\centering
\caption{API Integration Latency}
\label{tab:latency}
\begin{tabular}{llcc}
\toprule
\textbf{API} & \textbf{Protocol} & \textbf{Avg. Latency} & \textbf{Cache TTL} \\
\midrule
Open-Meteo & REST & 310\,ms & 300\,s \\
OpenSky Network & REST & 840\,ms & 60\,s \\
aisstream.io & WebSocket & 7.2\,s* & 60\,s \\
TomTom Flow & REST & 220\,ms & 300\,s \\
\bottomrule
\multicolumn{4}{l}{\footnotesize *Collection window; subsequent cached reads $<$\,10\,ms.}
\end{tabular}
\end{table}

The aisstream.io WebSocket latency reflects a 7-second collection
window designed to gather a statistically representative sample of
vessels in Indian waters. This cost is amortised over the 60-second
cache TTL; subsequent requests within this window return in under 10\,ms.

\subsection{Dashboard Performance}

End-to-end dashboard load time (home page, warm cache) averages
2.1\,seconds on a standard laptop browser over a 100\,Mbps connection.
Cold-cache load (all four APIs fetched simultaneously) averages
9.4\,seconds, primarily dominated by the aisstream.io collection
window. Streamlit's \texttt{@st.cache\_data} decorator ensures that
concurrent users share cached results, making cold-cache scenarios
rare in production.

\subsection{Cascade Failure Analysis}

Simulating failures of the top five hubs by betweenness centrality
reveals that removal of JNPT Port triggers the largest cascade,
isolating 3 secondary nodes and eliminating 14 trade route edges
(38.9\% of total edges). The post-failure risk delta averaged across
surviving edges is $+0.12$ (a 12-percentage-point increase in delay
probability), quantifying the systemic fragility introduced by heavy
reliance on a single seaport node.


% ─────────────────────────────────────────────────────────────────────────────
\section{Limitations and Future Work}

\subsection{Limitations}

\textbf{Class Imbalance:} The training dataset exhibits a 75:25
on-time to delayed ratio, causing the classifier to favour the
majority class. This suppresses recall on the critical delayed
class (0.253 for RF). Oversampling (SMOTE) and cost-sensitive
learning are planned for the next iteration.

\textbf{Graph Topology:} The current graph covers 15 nodes
representative of tier-1 Indian logistics infrastructure.
Tier-2 and tier-3 cities, which handle a significant proportion
of last-mile volume, are not yet modelled.

\textbf{aisstream.io Coverage:} AIS reception in Indian waters
is dependent on terrestrial receiver density. Remote coastal
areas may have sparse coverage, leaving a gap in vessel
tracking for offshore routes.

\textbf{TomTom Traffic:} The free tier (2,500 req/day) limits
road segment refresh to approximately one update per segment
every 15 minutes. High-frequency traffic events (accidents,
road closures) may not propagate immediately.

\subsection{Future Work}

\textbf{AfterShip Integration:} The \texttt{aftership\_service.py}
module is implemented and awaits an API key to ingest live
carrier-specific tracking events (Delhivery, Blue Dart, DTDC,
Xpressbees), which would provide ground-truth delay signals to
continuously retrain the model.

\textbf{NewsAPI Disruption Feed:} Integration of a news API
to automatically detect and alert on labour strikes, port
closures, and natural disasters would replace the current
heuristic alert thresholds with event-driven notifications.

\textbf{SMOTE and Rebalancing:} Applying SMOTE oversampling
to the minority class during training is expected to improve
F1 by approximately 10–15 percentage points based on pilot
experiments.

\textbf{Voice Interface:} A speech-to-text interface in regional
Indian languages would allow warehouse operators to query
shipment risk without navigating the web dashboard.


% ─────────────────────────────────────────────────────────────────────────────
\section{Conclusion}

This paper presented SupplyGuard, a real-time supply chain risk
intelligence platform that integrates four public APIs
(Open-Meteo, OpenSky Network, aisstream.io, and TomTom) into a
unified graph-theoretic risk model of Indian logistics infrastructure.
The system addresses a critical gap in existing SCRM tools by
providing continuously updated, multi-modal risk scoring that
reflects current weather, maritime, aviation, and traffic conditions
simultaneously.

Key contributions include: (1) a multi-source API fusion architecture
that updates route risk scores every 60 seconds; (2) SHAP-validated
evidence that port congestion, captured via live AIS data, is the
single most influential predictor of shipment delay; (3) a WCC-based
cascade failure simulator that quantifies second-order network
disruption; and (4) a production-ready Streamlit dashboard deployable
without cloud infrastructure.

Our experimental evaluation shows that the Random Forest
classifier outperforms Logistic Regression by F1 (0.333 vs.\ 0.216),
and that the end-to-end dashboard maintains sub-2.5-second refresh
latency under normal operating conditions. Future work will address
class imbalance, expand graph coverage, and integrate live shipment
tracking via AfterShip to close the feedback loop for continuous
model retraining.

% ─────────────────────────────────────────────────────────────────────────────
\begin{thebibliography}{00}

\bibitem{india_logistics_report}
FICCI--KPMG, ``Indian logistics sector: Opportunities and challenges,''
Federation of Indian Chambers of Commerce and Industry, Tech.\ Rep., 2022.

\bibitem{ecommerce_india}
Ministry of Commerce, ``E-commerce exports from India,''
Government of India, 2023.

\bibitem{chopra2004}
S. Chopra and M. S. Sodhi, ``Managing risk to avoid supply-chain
breakdown,'' \textit{MIT Sloan Management Review}, vol.~46, no.~1,
pp.~53--61, 2004.

\bibitem{kleindorfer2005}
P. R. Kleindorfer and G. H. Saad, ``Managing disruption risks in
supply chains,'' \textit{Production and Operations Management},
vol.~14, no.~1, pp.~53--68, 2005.

\bibitem{scor}
Supply Chain Council, ``Supply Chain Operations Reference Model
(SCOR) Version 12.0,'' 2017.

\bibitem{kup2022}
G. Kup, C. Dogu, and B. Dalmaç, ``Real-time prediction of delivery
delay in supply chains using machine learning approaches,''
\textit{Journal of Theoretical and Applied Information Technology},
vol.~100, no.~22, pp.~6612--6627, 2022.

\bibitem{jiang2019}
Y. Jiang, W. Shang, and Y. Liu, ``Predicting delays in cargo
transportation using gradient boosted trees on parcel data,''
\textit{IEEE Access}, vol.~7, pp.~167--178, 2019.

\bibitem{ivanov2016}
D. Ivanov, A. Dolgui, B. Sokolov, and M. Ivanova, ``Disruption-driven
supply chain (re)-planning and performance impact assessment with
consideration of pro-active and recovery policies,''
\textit{Transportation Research Part E}, vol.~90, pp.~7--24, 2016.

\bibitem{tang2008}
C. S. Tang and B. Tomlin, ``The power of flexibility for mitigating
supply chain risks,'' \textit{International Journal of Production
Economics}, vol.~116, no.~1, pp.~12--27, 2008.

\bibitem{ais_logistics}
S. Harati-Mokhtari, A. Wall, P. Brooks, and J. Wang,
``Automatic Identification System (AIS): A human factors approach,''
\textit{Journal of Navigation}, vol.~60, no.~3, pp.~373--389, 2007.

\bibitem{openmeteo}
Z. Zippenfenig, ``Open-Meteo: Weather Forecasts for Developers,''
\textit{Zenodo}, 2023. [Online]. Available:
\texttt{https://doi.org/10.5281/zenodo.7970649}

\bibitem{opensky}
M. Schäfer, M. Strohmeier, V. Lenders, I. Martinovic, and M. Wilhelm,
``Bringing up OpenSky: A large-scale ADS-B sensor network for
research,'' in \textit{Proc.\ IPSN}, 2014, pp.~83--94.

\bibitem{ecommerce_dataset}
M. Yagci, ``E-Commerce Shipping Dataset,''
\textit{Kaggle}, 2020. [Online]. Available:
\texttt{https://www.kaggle.com/datasets/prachi13/customer-churn}

\end{thebibliography}

\end{document}
